\section{Related Work}

\subsection{Long horizon RGB world models}

Driven by the abundance of data resources, numerous efforts to construct world models for autonomous driving have utilized RGB signals as their primary input to achieve robust long-horizon rollouts. For instance, GAIA-1 \cite{hu2023gaia} employs a multi-modal transformer architecture with language and action conditioning to autoregressively generate minute-long video clips. GAIA-2 \cite{russell2025gaia} extended this capability to multi-view RGB generation by transitioning the underlying backbone from a deterministic model to flow matching. Furthermore, Cosmos-Drive-Dreams \cite{ren2025cosmos} finetuning the Cosmos \cite{agarwal2025cosmos} foundation model to enhance the roll out of autonomous driving scene generation. In the indoor navigation task, Ego2Memory \cite{zhang2025mem2ego} explicitly uses VLM to map 2D features with Frontier Map for long-term persistent rollout.

By introducing a two-stage training and fine-tuning pipeline alongside better utilization of dynamic priors, Vista \cite{gao2024vista} significantly improves the resolution and realism of generated videos, as evidenced by FID and FVD metrics. Subsequently, DrivingWorld \cite{hu2024drivingworld} adopted a purely auto-regressive (AR) based approach. More recently, Epona \cite{zhang2025epona} provided the first open-source implementation for long video rollouts, utilizing an AR diffusion method.

\subsection{Occupancy World Models in Autonomous Driving}

While RGB images offer rich semantic features, accurately capturing and predicting the complex 3D geometry and physical depth of dynamic environments remains a significant challenge. Consequently, there is a growing trend towards developing world models based on 3D occupancy representations, which are typically derived from voxelized and densified LiDAR data.

Earlier works such as OccWorld \cite{zheng2024occworld} and Occllama \cite{wei2024occllama} addressed this by combining discrete encoding techniques (e.g., VQ-VAE \cite{van2017neural}) with deterministic auto-regressive (AR) transformer architectures, with Occllama further incorporating Large Language Models (LLMs) to enhance scene comprehension. Currently, generative modeling approaches have gained significant traction. For instance, DOME \cite{gu2024dome} utilizes a continuous VAE alongside a high-performance diffusion model, enabling fine-grained controllability. Recent advancements have further extended these capabilities: COME \cite{shi2025come} introduces a post-trained ControlNet \cite{zhang2023adding} approach, DynamicCity \cite{bian2024dynamiccity} leverages hex-plane decomposition, and OccFM \cite{liu2025towards} employs a Flow Matching-based multi-scale DiT \cite{peebles2023scalable} architecture. Expanding the scope, UniScene \cite{li2025uniscene} explores a conditional pipeline that uses BEV layouts to generate comprehensive multi-sensor data, ranging from RGB images to LiDAR point clouds. 

\subsection{Data driven simulation with generative models}

High-fidelity simulation is crucial for bridging the sim-to-real gap in autonomous driving, yet conventional physics-based simulators like CARLA \cite{Dosovitskiy17} often struggle with domain discrepancies. Recent research has shifted towards data-driven approaches to leverage large-scale real-world datasets. Nocturne \cite{vinitsky2022nocturne} facilitates high-efficiency simulation by replaying trajectories from the Waymo dataset. Building upon this, SLEDGE \cite{chitta2024sledge} and Scenario Dreamer\cite{rowe2025scenario} extends simulation to synthesized scenarios by generatively modeling road layouts and vehicle agents. In parallel, efforts have been made to generate more realistic sensory observations. ReSim \cite{yang2025resim} and DriveArena \cite{yang2025drivearena} utilize video generation techniques and HD maps to render high-resolution RGB observations for closed-loop policy evaluation. 

The raising of 3D Gaussian Splatting rendering in unbounded scene \cite{chen2024lidar, zhou2024hugsim, xiao2025splatco} allow us 
reconstruct static backgrounds more realistic. While, the SemCity\cite{lee2024semcity} and InifiCube\cite{lu2025infinicube} follow the out-painting solution introduced by Repainting\cite{lugmayr2022repaint} while combine it with tri-plane decomposition diffusion. Most recently, X-scene \cite{yang2025x} advances the field by proposing an external diffusion module that generates BEV layouts from natural language prompts, thus mitigating dependency on recorded driving logs. 
